\bibitem{HR}
D. Huybrechts and J. V. Rennemo,
\emph{Hochschild cohomology versus the Jacobian ring, and the Torelli theorem for cubic fourfolds},
Algebraic Geometry \textbf{6} (2019), 76--99.
\href{https://doi.org/10.14231/AG-2019-005}{doi:10.14231/AG-2019-005}.
Sections~2--3。一般のrotation付きHochschild環とJacobian環の比較の背景として参照する。
本稿の低次数の同定には次の\cite{LZ}を用いる。

\bibitem{LZ}
X. Lin and S. Zhang,
\emph{Serre algebra, matrix factorization and categorical Torelli theorem for hypersurfaces},
\href{https://arxiv.org/abs/2310.09927v2}{arXiv:2310.09927v2} (2024).
Theorem~2.2（dg比較とshift）、Theorem~3.5、Remark~3.6（sectorと積）、Section~5。
Proposition~3.7の互いに素という仮定はSerre代数への表示のための条件と区別する。

\bibitem{ExplicitUedaYoshinaga}
K. Ueda and M. Yoshinaga,
\emph{Logarithmic vector fields along smooth divisors in projective spaces},
Hokkaido Math. J. \textbf{38} (2009), 409--415.
\href{https://doi.org/10.14492/hokmj/1258553970}{doi:10.14492/hokmj/1258553970};
\href{https://arxiv.org/abs/0802.2133}{arXiv:0802.2133}.
Lemma~3：勾配イデアルとSebastiani--Thom型の関係。

\bibitem{ExplicitHwang}
J.-M. Hwang,
\emph{Symmetrizer group of a projective hypersurface},
J. Math. Soc. Japan \textbf{78} (2026), 55--62.
\href{https://doi.org/10.2969/jmsj/93469346}{doi:10.2969/jmsj/93469346};
\href{https://arxiv.org/abs/2501.08633}{arXiv:2501.08633}.
Theorems~1.5--1.6（勾配空間のfiberとsymmetrizer群）、
Proposition~2.2（可換代数の構造）。
